\documentclass{beamer}
\usepackage{xcolor}
\usepackage{makecell}
\usepackage{array}
\usepackage{ragged2e}
\newcolumntype{P}[1]{>{\RaggedRight\hspace{0pt}}p{#1}}
\begin{document}
\title{RuleBreaker: Categorical Correlations as Probabilistic Rules}   
\author{Nathan Danneman} 
\date{\today} 

\logo{%
    \includegraphics[width=1.5cm,height=1cm,keepaspectratio]{dm_logo}
}

\frame{\titlepage} 

%\-\hspace{1cm}  indented text

\frame{
\frametitle{Acknowledgements}
Big thanks to:
\begin{itemize}
\item BSides Asheville
\item Data Machines (www.datamachines.io)
\item Uncle Sam
\end{itemize}
\bigskip \bigskip \bigskip
 \textcolor{gray}{This research was developed with funding from the Defense Advanced Research Projects Agency (DARPA).} \\
 \bigskip
 \textcolor{gray}{The views, opinions and/or findings expressed are those of the author and should not be interpreted as representing the official views or policies of the Department of Defense or the U.S. Government.}

}


\frame{
\frametitle{Outline}
\begin{enumerate} 
\item ``Big data'' and analytics for cyber defense \\
\item RuleBreaker: how it works \\
\item RuleBreaker: security relevance, use cases \\
\item Next steps, etc.
\end{enumerate}
}

%
%\frame{
%\frametitle{Defensive Analytics I: Rules}
%Idea: Encode knowledge about known TTPs as rules.\\
%\bigskip
%\begin{itemize}
%\item e.g. Snort \\
%\item If (ip AND port OR domain) then (block/report/etc) \\
%\pause
%\item Pros: understandable, targeted, fast (to implement and execute)\\
%\item Cons: simplistic; known-knowns; short shelf-life
%\end{itemize}
%}
%
%
%\frame{
%\frametitle{Defensive Analytics II: Statistical Pattern Matching}
%Idea: Train a model to differentiate malicious cases from safe ones.\\
%\bigskip
%\begin{itemize}
%\item e.g. FireEye \\
%\item Statistical model separating, e.g., malware-containing .doc(x) from safe examples
%\pause
%\item Pros: targeted, accurate \\
%\item Cons: need training data; uninterpretable models; short shelf-life; need many models to achieve broad coverage
%\end{itemize}
%}
%
%\frame{
%\frametitle{Defensive Analytics III: Statistical Anomaly Detection}
%Idea: Use a statistical model to highlight anomalous activity.\\
%\bigskip
%\begin{itemize}
%\item e.g. IronNet \\
%\item Either model normalcy and alert on deviations, or directly estimate ``anomalousness''
%\pause
%\item Pros: dynamic models; per-environment customization; detect unknowns  \\
%\item Cons: noisy; definition of ``anomaly'' implies algo implies variance in results; uninterpretable
%\end{itemize}
%}
%
%
%\frame{
%\frametitle{Best of All Worlds?}
%RuleBreaker attempts to get some of the power of each approach, while avoiding each one's pitfalls.\\
%\bigskip
%
%Idea: generate simplified models of normalcy in the form of probabilistic rules.
%
%\bigskip
%\begin{itemize}
%\item Dynamism and `unknown unknowns' of anomaly detection \\
%\item Interpretability of rules \\
%\item Accuracy of statistical pattern matching
%\end{itemize}
%
%}
%
%
%\frame{
%\frametitle{Motivation}
%Network data, through RFCs and convergent usage patterns, often includes strong interdependencies between fields.
%\bigskip
%\begin{itemize} 
%\item File extensions imply mime\_type.\\
%\item Port 53 traffic + AAAA record implies range of return bytes.\\
%\item Port number correlates strongly with protocol.
%\end{itemize}
%\bigskip
%Idea: Automatically identify strong, categorical correlations to \textcolor{red}{learn about a network}. \\
%\bigskip
%Then, find observations that break them, to \textcolor{red}{identify anomalies}.
%}
%
%\frame{
%\frametitle{Further Motivation}
%
% \begin{overlayarea}{\textwidth}{0.23\textheight}
% \visible<1>{
%\includegraphics[width=1.9cm]{thinking_face}
%\end{overlayarea}
%}
%\visible<1->{
%Can't we just use SQL to find these?
%}
%\visible<2->{
%\begin{itemize}
%\item Yes, \emph{if you already know} what mis-matches you're looking for
%\item Yes, \emph{if there are few mis-matches} you care about
%\end{itemize}
%}
%
%\bigskip
%\visible<3->{
%Hey, isn't this just anomaly detection?
%\begin{itemize}
%\item Yes-ish.
%\item Optimized for categorical data
%\item Performant in the face of irrelevant variables
%\item Focused on `near-misses'
%\end{itemize}
%}
%}
%
%
%\frame{
%\frametitle{The Language of Rules}
%RuleBreaker thinks of categorical correlations as \emph{rules}.\\
%\bigskip
%Example: if DNS, then port 53 with probability 98\%.\\
%\bigskip
%These are \emph{descriptive} rather than \emph{normative}; they are identified from data.\\
%\bigskip
%Rules have three parts: predicate, consequent, and (conditional) likelihood.
%
%}
%
%
%\frame{\frametitle{Inducing Rules from Data}
%Association Rule Mining is the task of identifying antecedent sets that are highly likely to imply consequents.\\ 
%\textcolor{gray}{(ie. find strong correlations in categorical data)}\\
%\bigskip
%Originated with market basked analysis\\
%\textcolor{gray}{(bread + peanut butter $=>$ jelly)}\\
%\textcolor{gray}{see Agrawal, Imielinski, and Swami (1993)}\\
%\bigskip
%Task involves first identifying frequent itemsets \\
%\textcolor{gray}{(bread, peanut butter, jelly)}\\
%\bigskip
%Next, you use those itemsets to identify high-confidence rules\\
%\textcolor{gray}{ \{bread, peanut butter $=>$ jelly 99\% of the time \} }
%}
%
%\frame{\frametitle{Identifying Frequent Itemsets}
%
%A \textcolor{red}{Frequent Itemset} is a set (group of levels of variables) that occurs with higher than S(upport) frequency in a database.\\
%\bigskip
%FP-Growth \textcolor{gray}{(Han et al, 2000)}
%\begin{itemize}
%\item Scan database ONCE to find (sorted list) of frequent items
%\item Grow a tree based on order of frequency
%\item Use the tree structure to rapidly identify frequent itemsets without candidate generation
%%\item Check for confidence of rules made from binary partitions of itemsets
%\end{itemize}
%}
%
%
%\frame{\frametitle{Identifying Rules from Frequent Itemsets}
%A \textcolor{red}{Rule} is a unary split of a set where the likelihood of the consequent given the predicate in the dataset has a minimum \emph{confidence} (percentage).\\
%\bigskip
%
%Itemset: \{A,B,C\}\\
%Unary split: \{A,C\}, \{B\}\\
%Rule: If \{A,C\} then \{B\} with probability 96\%\\
%\bigskip
%
%Rule Induction:
%Generate unary splits of all frequent itemsets, and check the resulting conditional probabilities in the database.
%}
%
%
%\frame{\frametitle{Why Do We Care?}
%
%Learn about \textcolor{red}{patterns} in the network:\\
%
%\begin{itemize}
%\item DNS running over an atypical port
%\item Services associated with username-types
%\item Note proxies
%\end{itemize}
%
%\bigskip
%
%Identify \textcolor{red}{discrepancies}.\\
%E.g. ``IF port=53 THEN protocol=`dns' 99.7\% of cases.''~~~~~(rule)\\
%1.2.3 sent icmp traffic to 4.5.6 over port 53.~~~~~~~~~~~~~~~~(breakage)
%
%}
%
%\frame{\frametitle{How to Generate Security-Relevant Rules}
%
%At a high level, RuleBreaker is anomaly detection. \\
%
%\bigskip
%
%How can we induce \textcolor{red}{relevant} rule?
%
%\begin{itemize}
%\item Choose the right features
%\item Make sure the rules are statistically sound
%\item Automate the dropping of uninteresting rules
%\end{itemize}
%
%}
%
%
%
%\begin{frame}
%\frametitle{Security-Relevant Rules and Breakages: Features}
%\begin{center}
%\begin{tabular}{ P{3.5cm} | P{7cm} }
% Features  &  Rationale \\ 
% \hline \hline
% 
% \makecell[l]{usernames\\source/dest ips} & Identify patterns and discrepancies in login locations \\
% \hline
% \makecell[l]{sourceIP\\destIP\\destPort} &Identify common connection patterns and anomalies (possible internal probing) \\
% \hline
%user agent tokens & Identify bogus UAs \\
% \hline
% \makecell[l]{(inferred) file type\\file extension} & File maqeurading \\
% \hline 
% \makecell[l]{destination port\\(inferred) protocol} & Atypical connections \\
% \hline 
% \makecell[l]{HTTP Method\\(inferred) file type} & e.g. ``POST'' to images, CSS \\
% \hline 
%
%\end{tabular}
%\end{center}
%
%\end{frame}
%
%
%\begin{frame}
%\frametitle{Identifying Security-Relevant Anomalies: Metrics}
%Not all rules are \textcolor{red}{statistically interesting}.\\
%\bigskip
%Consider the rule: If A, then B. \\
%\bigskip
%Confidence: $\frac{Support(A~and~B) }{Support(A)}$\\ 
%\medskip
%\textcolor{gray}{equivalently: pr(B$\vert$A)} \\
%\pause
%
%\bigskip
%Lift: $\frac{Support(A~and~B)}{Support(A) ~* ~Support(B)}$ \\
%\medskip
%\textcolor{gray}{penalizes common B's }
%
%\end{frame}
%
%
%\frame{
%\frametitle{Confidence vs Lift}
%\includegraphics[width=10cm]{venn}
%
%}
%
%
%\begin{frame}[fragile]
%\frametitle{Identifying Security-Relevant Anomalies: Rule Pruning}
%Not all breakages or all rules are \textcolor{red}{relevant}.\\
%\bigskip
%Rule: if mime\_type is text/html, extension is .html with probability 0.98 \\
%\medskip
%Breakage: a row with mime\_type text/html and extension .txt \\
%
%\bigskip
%We implemented an \emph{exceptions list} that allows rule-and-breakage sets to be omitted from findings list.  These specify the match criteria (Any, Equals, Contains, Regex), and the three portions of a finding (predicate, consequent, and breakage).\\
%\medskip
%
%Example:
%\begin{verbatim}
%a();;a();;c(ext:jpeg;mime:image)
%\end{verbatim}
%
%\end{frame}
%
%
%\begin{frame}[fragile]
%\frametitle{Example: Bro\_Files Use Case}
%We think inferred mime\_type should be tightly coupled with file extension (or vice versa?) \\
%
%\bigskip
%
%\begin{itemize}
%\item Anecdoatal evidence \\
%\item Unknown unknowns
%\end{itemize}
%\end{frame}
%
%
%\begin{frame}[fragile]
%\frametitle{Example:  Bro\_Files\_View Rules}
%Identified the following rules and conditional probabilities: \\
%\textcolor{gray}{Read as: ``antecedent implies consequent with confidence X}
%\begin{verbatim}
%[txt] => [text/plain], 0.961
%[xml] => [application/xml], 0.999
%[application/xml] => [xml], 1.0
%[application/vnd.ms-pol] => [pol], 1.0
%[pol] => [application/vnd.ms-pol], 0.989
%[application/x-dosexec] => [dll], 0.999
%[dll] => [application/x-dosexec], 1.0
%\end{verbatim}
%\end{frame}
%
%
%\begin{frame}
%\frametitle{Example:  Bro\_Files Findings}
%
%Rule: If application/x-dosexec, then extension '.dll' (99\%) \\
%\smallskip
%Breakage: A IP pulling down an x-dosexec filetype with a '.txt' file extension. \\
%\bigskip
%Note: This is a \textcolor{red}{starting point} for an investigation; rarely a standalone finding.
%\end{frame}
%
%
%\begin{frame}
%\frametitle{Where Can I Get One?}
%\begin{enumerate}
%\item Rules
%\begin{itemize} 
%\item Spark: frequent pattern mining module in MLlib 
%\item R: arules package
%\item Python: mlxtend.frequent\_patterns
%\end{itemize}
%\item Rule Checker: Write your own
%\item Most Complex Rule: Write your own
%\item Confidence baked in; R/Python have lift
%\item Exception List: You are going to want one
%\end{enumerate}
%\end{frame}
%
%
%\frame{\frametitle{Stengths and Limitations}
%Strengths:
%\begin{itemize}
%\item Broad applicability
%\item Explainability
%\item Speed and scalability
%\item Tuning is trivial
%\end{itemize}
%\pause \medskip
%Weaknesses:
%\begin{itemize}
%\item High confidence rules are not automatically interesting
%\item Categorical data/binning
%\item Only identifies ``orthogonal outliers''
%\item Does not handle multi-item consequents
%\end{itemize}
%}


% [ext:.json] => [mime:text/json]                | ext:.json;mime:application/x-dosexec 

\frame{
\frametitle{}
IF \textcolor{red}{lunch time} THEN \textcolor{red}{not too many questions}?\\
\bigskip
\bigskip
nathandanneman [at] datamachines [dot] io
}


%
%\frame{\frametitle{Slide Title}
%Slide Contents
%}
%
%\frame{\frametitle{Slide Title}
%Slide Contents
%}
%
%\frame{\frametitle{Slide Title}
%Slide Contents
%}





%
%\section{Section no.1} 
%\frame{\frametitle{Title} 
%Each frame should have a title.
%}
%\subsection{Subsection no.1.1  }
%\frame{ 
%Without title somethink is missing. 
%}
%
%
%\section{Section no. 2} 
%\subsection{Lists I}
%\frame{\frametitle{unnumbered lists}
%\begin{itemize}
%\item Introduction to  \LaTeX  
%\item Course 2 
%\item Termpapers and presentations with \LaTeX 
%\item Beamer class
%\end{itemize} 
%}
%
%\frame{\frametitle{lists with pause}
%\begin{itemize}
%\item Introduction to  \LaTeX \pause 
%\item Course 2 \pause 
%\item Termpapers and presentations with \LaTeX \pause 
%\item Beamer class
%\end{itemize} 
%}
%
%\subsection{Lists II}
%\frame{\frametitle{numbered lists}
%\begin{enumerate}
%\item Introduction to  \LaTeX  
%\item Course 2 
%\item Termpapers and presentations with \LaTeX 
%\item Beamer class
%\end{enumerate}
%}
%\frame{\frametitle{numbered lists with pause}
%\begin{enumerate}
%\item Introduction to  \LaTeX \pause 
%\item Course 2 \pause 
%\item Termpapers and presentations with \LaTeX \pause 
%\item Beamer class
%\end{enumerate}
%}
%
%\section{Section no.3} 
%\subsection{Tables}
%\frame{\frametitle{Tables}
%\begin{tabular}{|c|c|c|}
%\hline
%\textbf{Date} & \textbf{Instructor} & \textbf{Title} \\
%\hline
%WS 04/05 & Sascha Frank & First steps with  \LaTeX  \\
%\hline
%SS 05 & Sascha Frank & \LaTeX \ Course serial \\
%\hline
%\end{tabular}}
%
%
%\frame{\frametitle{Tables with pause}
%\begin{tabular}{c c c}
%A & B & C \\ 
%\pause 
%1 & 2 & 3 \\  
%\pause 
%A & B & C \\ 
%\end{tabular} }
%
%
%\section{Section no. 4}
%\subsection{blocs}
%\frame{\frametitle{blocs}
%
%\begin{block}{title of the bloc}
%bloc text
%\end{block}
%
%\begin{exampleblock}{title of the bloc}
%bloc text
%\end{exampleblock}
%
%
%\begin{alertblock}{title of the bloc}
%bloc text
%\end{alertblock}
%}
\end{document}